%%%%%%%%%%%%%%%%%%%%%%%%%%%%%%%%%%%%%%%%%%%%%%%%%%%%%%%%%%%%%%%%%%%%%
\part{The Execution Layer}\label{part:execution}
%%%%%%%%%%%%%%%%%%%%%%%%%%%%%%%%%%%%%%%%%%%%%%%%%%%%%%%%%%%%%%%%%%%%%

Part~\ref{part:protocol} described infrastructure that proves that a given asset state is spent exactly once. This guarantee supports most blockchain applications, including issued currencies, non-fungible assets, programmable ownership, automated market making, and bridges to external chains. These applications do not require a globally shared state machine.

\section{Tokens}

A token is a self-contained ledger containing its genesis, complete transfer sequence, and certificates proving that each transfer was recorded exactly once. A recipient establishes validity by checking these records against the trust base, without consulting external state. A wallet can therefore run in a browser tab, and verification can be performed offline.

\begin{figure}[htbp]
    \centering
    \begin{tikzpicture}[diagram,
        tx/.style={diagram box, minimum width=2.75cm, minimum height=1.25cm, align=center,
                   font=\scriptsize},
        cert/.style={diagram box faint, minimum width=2.15cm, minimum height=0.5cm,
                     font=\scriptsize}]
        \node[tx, fill=uniFillB, draw=uniAccent] (g) at (0,0)
            {\textbf{Genesis}\\token type\\mint reason\\first owner};
        \node[tx] (t1) at (3.55,0) {\textbf{Transfer}\\new owner\\condition};
        \node[tx] (t2) at (7.1,0) {\textbf{Transfer}\\new owner\\condition};
        \node[tx] (t3) at (10.65,0) {\textbf{Transfer}\\current\\owner};

        \foreach \a/\b in {g/t1, t1/t2, t2/t3}{
            \draw[diagram arrow] (\a) -- (\b);
        }

        \foreach \n in {g,t1,t2,t3}{
            \node[cert] (c\n) at ($(\n)+(0,-1.35)$) {certificate};
            \draw[diagram connector, draw=uniLine!55] (\n) -- (c\n);
        }

        \node[diagram box emph, minimum width=2.5cm, align=center] (r) at (13.6,0)
            {Recipient};
        \draw[diagram arrow] (t3) -- (r);

        \draw[diagram brace] (-1.5,-2.15) -- (12.15,-2.15);
        \node[diagram caption, align=center] at (5.3,-2.6)
            {verified end to end by the recipient, against the trust base alone};
    \end{tikzpicture}
    \caption{A self-contained token ledger. Each step includes its recording certificate. The recipient re-verifies the complete chain without consulting external state.}
    \label{fig:token}
\end{figure}

\subsection{States and Identifiers}

A token state is a spending condition together with a hash linking it to the token's history. The \emph{state identifier} is derived from that pair, and it is the key the aggregation layer records when the state is spent. Each transfer introduces fresh randomness into the next state hash, so the next state identifier is unpredictable to anyone who has not seen the transaction. The aggregation layer therefore cannot connect one spend of a token to the next, even though it records both.

\medskip

The network observes a sequence of apparently unrelated identifiers. It receives no amounts, asset types, owner details, or links between spends of the same token.

\subsection{Minting and Issuance Authority}

The minting key pair is a public constant of the protocol. Anyone can construct a syntactically valid mint transaction. A certified mint proves only that the token identifier was fresh; it confers no issuance authority.

\medskip

Issuance authority is established by two genesis fields that every subsequent recipient rechecks.

\medskip

The token type classifies the token and specifies its permanent rules. A type definition fixes payload decoding, admissible minting justifications, and issuance policy. A published definition is immutable. Every verifier must apply byte-identical rules indefinitely to prevent different recipients from reaching different validity decisions for the same token.

\medskip

The mint reason justifies the mint under the type's rules. It can be an issuer's authorization signature, a value-conservation proof for change from a split, a proof of a matching asset locked on an external chain, or a zero-knowledge proof of any statement accepted by the type.

\medskip

Recipients enforce issuance rules during verification. Validators may include a counterfeit mint, but recipients reject it because it fails the required checks. This fail-shut behavior is mandatory on the public network. A verifier must reject a token if its type is unknown, its reason format is unregistered, or its reason is inadmissible for its type. Such tokens cannot be accepted as unvalidated.

\subsection{Issuance Mechanisms}

\begin{description}
    \item[Native issuance.] An issuer defines a type pinning its own key and authorizes each mint with a signature. This is the ordinary case for a currency or a collection.

    \item[Derived issuance.] A token can represent change from a split. A split burns a source token and mints outputs whose combined value provably equals the source value. A merge performs the reverse operation. In both cases the reason is a conservation proof governed by a network-wide rule, independent of the issuer.

    \item[Bridged issuance.] The reason is a proof that a matching amount is locked in a vault on another chain, checked by every recipient against that chain's own consensus. Section~\ref{sec:bridging} covers this in full.

    \item[Proof-carrying issuance.] The reason is a succinct proof of an arbitrary statement defined by the type. Any statement that a verifier can check in constant time can authorize a mint.
\end{description}

\subsection{Value, Splitting, and Conservation}

A fungible token carries a collection of assets, each an identifier and a non-negative amount. Splitting divides it while preserving every enclosed asset independently.

\medskip

The construction is a sum tree. The owner allocates the source value across the intended outputs, builds one tree per asset identifier whose leaves are the outputs and whose root commits to the total, and then burns the source token in a transfer whose data is exactly that commitment. Each output is then minted carrying the burned source token and an inclusion proof against the committed allocation.

\medskip

A verifier of any output re-derives its commitment from the mint fields and verifies inclusion in the allocation tree. It also requires the tree's committed total to equal the source token's amount. The recipient checks value conservation independently of the wallet that performed the split. A recipient learns the source amount, verifies the source token completely, and accepts only if no value could have been created.

\medskip

Across the platform, negligence can destroy value but cannot create it. Burning a source token and then losing the data needed to mint an output destroys that output's value. No party can inflate the supply.

\subsection{Token Standards}

Type definitions are published in a registry, and a set of standard type templates provides interoperability.

\begin{table}[ht]
    \centering
    \caption{Standard token type templates}
    \label{tab:uts}
    \begin{tabular}{@{}llp{7.7cm}@{}}
        \toprule
        \textbf{Standard} & \textbf{Shape} & \textbf{Notes} \\
        \midrule
        UTS-0 & Blank container & Unrestricted minting, application-defined payload. Proves only that someone minted it. Wallets must not present these as carrying value. \\
        \addlinespace
        UTS-1 & Issued fungible asset & Analogous to ERC-20. Issuer-gated primary issuance, divisible by split, optionally consolidatable by merge. \\
        \addlinespace
        UTS-2 & Non-fungible token & Analogous to ERC-721. The payload commits to immutable content by hash. Curated collections gate minting by issuer; open collections do not. \\
        \addlinespace
        UTS-3 & Capped-issuance series & A supply cap with trustless enforcement. Every valid issuance consumes one of a bounded set of genesis states, so no party needs to observe all issuance to know the cap held. \\
        \addlinespace
        UTS-4 & Pooled assets & Value aggregated under type-defined rules. See Section~\ref{sec:pools}. \\
        \bottomrule
    \end{tabular}
\end{table}

UTS-3 provides a locally verifiable supply cap. On a shared ledger, cap verification depends on the correctness of the enforcing contract and its execution by the chain. In UTS-3, the type definition specifies the cap. Each issuance must consume one member of a bounded set of genesis states, each admitted at most once by the aggregation layer. A recipient verifies compliance with the cap for the received token without observing other issuance. The cap applies per network; deploying the same definition on a second network does not create a combined cap across the two.

\section{Transaction Flow and Latency}

A transfer involves three parties and one round trip to the network. The recipient supplies the spending condition under which it wishes to receive the token. The sender builds the transaction, satisfies the current condition, and asks the aggregation layer to record the state being spent. The layer returns a proof. The sender passes the token, now extended by one transfer and its proof, to the recipient, who verifies it.

\medskip

Settlement is final when the proof is issued, in one to two seconds. For workloads requiring lower latency, the protocol supports a second ordering of the same steps.

\begin{figure}[htbp]
    \centering
    \begin{tikzpicture}[diagram,
        stepbox/.style={diagram box, minimum width=2.45cm, minimum height=1.3cm,
                     align=center, font=\scriptsize},
        donebox/.style={diagram box emph, minimum width=2.45cm, minimum height=1.3cm,
                     align=center, font=\scriptsize}]

        % ---------- Standard path ----------
        \node[diagram label plain, font=\small\bfseries, anchor=west] at (-6.95,1.35)
            {Standard path};
        \node[stepbox] (a1) at (-5.72,0) {Sender builds\\the transfer};
        \node[stepbox, fill=uniFillB, draw=uniAccent] (a2) at (-2.86,0)
            {Aggregation\\records the state};
        \node[stepbox] (a3) at (0,0) {Sender passes\\token and proof};
        \node[stepbox] (a4) at (2.86,0) {Recipient\\verifies};
        \node[donebox] (a5) at (5.72,0) {Goods\\released};
        \foreach \x/\y in {a1/a2, a2/a3, a3/a4, a4/a5}{\draw[diagram arrow] (\x) -- (\y);}
        \draw[diagram brace] (-4.15,-0.82) -- (-1.57,-0.82);
        \node[diagram caption, align=center] at (-2.86,-1.3)
            {settled here: one to two seconds};

        % ---------- Low-latency path ----------
        \node[diagram label plain, font=\small\bfseries, anchor=west] at (-6.95,-2.35)
            {Low-latency path};
        \node[stepbox] (b1) at (-5.72,-3.7) {Sender passes\\token and transfer};
        \node[stepbox] (b2) at (-2.86,-3.7) {Recipient submits\\to aggregation};
        \node[stepbox, fill=uniFillB, draw=uniAccent] (b3) at (0,-3.7)
            {Non-inclusion\\returns at once};
        \node[donebox] (b4) at (2.86,-3.7) {Goods\\released};
        \node[stepbox] (b5) at (5.72,-3.7) {Inclusion follows:\\token spendable};
        \foreach \x/\y in {b1/b2, b2/b3, b3/b4, b4/b5}{\draw[diagram arrow] (\x) -- (\y);}
        \draw[diagram brace] (-1.28,-4.52) -- (1.28,-4.52);
        \node[diagram caption, align=center] at (0,-5.0)
            {goods released before the round completes};
    \end{tikzpicture}
    \caption{Two orderings of the same transfer. In the low-latency path the recipient submits the request itself and acts on the immediate confirmation that the state had not been spent, waiting for the inclusion proof only before treating the received token as spendable inventory.}
    \label{fig:latency}
\end{figure}

In the low-latency path, the sender delivers the token and transaction directly to the recipient, who submits the request. The layer immediately confirms that the state has not been spent and is scheduled for recording. The recipient can then release goods: only the sender could have spent the state previously, and the recipient now controls the request. The inclusion proof follows in the next round and is required before the token is treated as spendable inventory. With sufficient inventory to cover the round time, a participant can transact continuously without waiting for each round.

\medskip

Neither path involves a mempool, a block, a fee auction, or a confirmation depth. A party that is offline can verify a token it received an hour ago or a year ago with the same procedure and the same result.

\section{Programmable Ownership}

Ownership in Unicity can be specified by a \emph{predicate}, an arbitrary rule determining when a token may be transferred. Predicates generalize public-key ownership and are evaluated off-chain by the transaction participants.

\medskip

The standard set covers the usual requirements: a signature under a given key, payment to a key hash, payment to a script hash, multi-signature requiring all of a set, threshold signature requiring some of a set, a timelock that forbids transfer before a reference time, a hashed timelock combining a secret with a deadline, a burn predicate that can never be unlocked, and an open predicate that anyone may satisfy. Time-dependent predicates read their reference time from the certified round rather than from any local clock, so two verifiers always agree about whether a deadline had passed.

\medskip

A reduction to unforgeability of the predicate family proves that generalizing ownership from public keys to predicates preserves three execution-layer properties. A state cannot be spent twice, only the legitimate owner can prevent a token from being spent, and the network cannot link transactions of the same token.

\subsection{Atomic Swaps}

Two parties can exchange tokens without trusting each other or any intermediary, using interlocking predicates that guarantee either both transfers complete or neither does. Correctness is enforced by the predicates themselves. There is no custodian, no escrow contract holding both sides, and no global state in which the swap is arranged.

\medskip

For applications requiring agreement among several parties, the parties coordinate directly. Conditions attached to their tokens enforce the agreed outcome.

\section{Verifiable Execution}\label{sec:verifiable}

Execution is \emph{verifiable} when acceptance is a deterministic function of its output and a rule set fixed in advance and identical for every party. A recipient checks the result without relying on the producer.

Verification requires immutable rules to prevent disagreement between verifiers. Checks must be inexpensive enough for recipients to perform, and failure behavior must require rejection of unverifiable tokens rather than provisional acceptance. Verification resources must also be bounded so that malicious inputs exhaust a defined budget and are rejected without exhausting the verifier.

\medskip

Unicity provides three mechanisms for verifiable execution with different costs.

\subsection{Re-Verification Against a Frozen Rule}

By default, every recipient reruns the token type's rules to check the producer's computation. This enforces issuance policies, value conservation, and type-defined state machines without consensus.

\medskip

Verification work grows with the token's provenance. A worklist initially contains the received token. If verification reveals a genesis dependency on an earlier token, that token is added to the worklist; verification continues until the list is empty. A single cumulative budget limits bytes, items, nesting depth, and proof steps across the entire run. A deliberately deep or wide provenance graph that exhausts the budget is rejected as invalid. Identically configured verifiers accept exactly the same tokens.

\subsection{Succinct Proofs}

A succinct proof can replace expensive re-verification. The producer generates an argument of correct computation that the verifier checks at a cost independent of the computation's size. This mechanism supports split conservation over large allocations, token history compression, and bridge redemption. In bridge redemption, a contract on another chain verifies the proof because replaying the underlying history would be too expensive.

\subsection{Shared, Ordered State}

Some applications require multiple parties to write to one object in a defined order, such as an order book with a single matching queue. The enshrined EVM partition provides a conventional execution environment and guarantees for these applications, certified by the same consensus layer. Section~\ref{sec:evm} describes this partition.

\subsection{Execution That Cannot Be Re-Checked}

Not all useful computation is deterministic or can be recomputed by a recipient. Model inference, for example, cannot be cheaply recomputed and often cannot be recomputed at all. Assurance mechanisms for such computation introduce trust assumptions beyond those of the cryptographic core.

\begin{table}[h]
    \centering
    \caption{Assurance mechanisms for computation that recipients cannot rerun}
    \label{tab:verifexec}
    \begin{tabular}{@{}lp{5.4cm}p{4.3cm}@{}}
        \toprule
        \textbf{Mechanism} & \textbf{How it works} & \textbf{What must be trusted} \\
        \midrule
        Trusted execution & Hardware attests that specific code ran in an enclave & The hardware manufacturer \\
        \addlinespace
        Fraud proofs & Assume correct, allow a challenge within a window & At least one honest watcher, plus the delay \\
        \addlinespace
        Attestation committee & A set of parties independently execute and attest & An honest majority of that set \\
        \bottomrule
    \end{tabular}
\end{table}

Unicity supports the use of these assurance mechanisms. An agent performing inference can hold tokens, receive payments, and settle on the platform. Its payments retain the full cryptographic guarantees regardless of how the inference was produced. These guarantees apply only to settlement; the model's output is not cryptographically verified.

\section{Autonomous Pools}\label{sec:pools}

An autonomous pool implements an automated market maker without an operator.

\medskip

A pool is a single token whose state contains two asset reserves, the total shares outstanding, and an accumulator recording every processed deposit and withdrawal. Its owner predicate is the open predicate, allowing anyone to spend the current state. The type's rules permit only transitions satisfying the pool's equations and require the successor to retain the open predicate. No key grants administrative privileges, so no party can take private control of the pool.

\begin{figure}[htbp]
    \centering
    \begin{tikzpicture}[diagram,
        act/.style={diagram box, minimum width=2.6cm, minimum height=0.95cm,
                    align=center, font=\scriptsize}]
        \node[diagram box emph, text width=4.3cm, align=center, minimum height=1.9cm]
            (st) at (0,0)
            {\textbf{Pool state}\\[2pt]
             {\scriptsize reserves of two assets}\\
             {\scriptsize shares outstanding}\\
             {\scriptsize accumulator of consumed burns}\\[3pt]
             {\scriptsize owner: \emph{open predicate}}};

        \node[act] (dep) at (-5.0,1.7) {\textbf{Deposit}\\burn assets};
        \node[act] (swp) at (-5.0,0.0) {\textbf{Swap}\\burn input asset};
        \node[act] (wd)  at (-5.0,-1.7) {\textbf{Withdraw}\\burn shares};

        \foreach \n in {dep,swp,wd}{
            \draw[diagram arrow] (\n.east) -- (st.west);
        }

        \node[act] (osh) at (5.2,1.7) {shares\\minted};
        \node[act] (oas) at (5.2,0.0) {output asset\\minted};
        \node[act] (ow)  at (5.2,-1.7) {reserves\\returned};
        \foreach \n in {osh,oas,ow}{
            \draw[diagram arrow up] (st.east) -- (\n.west);
        }

        \node[diagram caption, text width=13.6cm, align=center] at (0,-3.0)
            {Every transition is re-checked by every verifier of every token the pool produces.
             The party advancing the state bears the cost and receives no gain for doing so.};
    \end{tikzpicture}
    \caption{An autonomous pool. Value enters by burning a token bound to the pool, and leaves as a commitment the claimant mints against. No party holds custody at any point.}
    \label{fig:pool}
\end{figure}

A depositor burns a token committing to the pool and the intended share recipient. A swapper burns an input asset committing to a minimum acceptable output and a refund recipient. A withdrawer burns shares. Each burn identifies the pool without specifying a particular pool state. A party advancing the pool consumes a burn, applies the corresponding equation, inserts the burn into the accumulator to prevent repeated consumption, and commits to the output the claimant will mint.

\medskip

The pool rules enforce the following properties.

\begin{description}
    \item[Solvency.] Every unit of value in the reserves entered through a certified burn counted exactly once, and every exit decrements the reserves by construction. The pool cannot commit to releasing more than it holds.

    \item[Conservation across types.] The total supply of the underlying asset, including ordinary tokens and pool reserves, is invariant. Every verifier of a token minted by a withdrawal re-verifies the pool's history through the same worklist used for other provenance.

    \item[No dilution.] Value per share is non-decreasing under deposits and withdrawals, and the constant-product rule with a fee means the reserve product never decreases across a swap. No sequence of trades, at any prices, pushes the value backing a share below the curve.

    \item[No capture.] Every transition must retain the open predicate and satisfy the state equations. No code path allows any party, including the pool's creator, to extract value except by burning shares.

    \item[Race safety.] When two parties attempt to advance the same pool state, exactly one succeeds. The other party's burn remains available for a later transition because it commits to the pool rather than a particular state.

    \item[No stranding.] A refund transition can consume a burn whose bound is unsatisfiable, such as a swap whose minimum output can no longer be met. It returns the burn's exact contents to the committed refund recipient. A hostile party that advances the pool can decline a trade but cannot confiscate its value.
\end{description}

A pool has two limitations. Its single spendable state serializes operations and creates contention. The party that successfully advances the pool also chooses the order of applied pending operations, permitting reordering as in other automated market makers. The minimum-output bound limits losses on each trade, and the refund path limits the outcome to cancellation when the bound cannot be satisfied. This exposure arises from the application's shared pool object.

\medskip

Other applications can use the same mechanism. Lending can lock collateral under conditions that release it on liquidation terms. Auctions can hold bids under conditions tied to winner selection, and games can hold stakes under conditions tied to outcomes. Ordinary software coordinates these operations, while rules rerun by every recipient enforce them.

\section{Privacy}

Unicity limits the data submitted to the network to a state identifier and the data satisfying a spending condition. The amount, asset, sender, recipient, and token are never submitted. Their privacy therefore does not depend on encrypting or hiding submitted data.

\medskip

Two consecutive spends of the same token are unlinkable to the network, because the state mask makes the next identifier unpredictable. What remains is that a party holding the token can see its history, and that a party who transacts repeatedly under one key can be recognized by counterparties who see that key.

\medskip

The remedy is a fresh spending condition per transaction, which would ordinarily mean managing a fresh keypair per transaction. Unicity provides protocols that derive unlinkable conditions from a single persistent keypair, in interactive and non-interactive forms, the latter allowing a published address to receive payments that share no visible key. These operate at the wallet level and require no support from the protocol, so a user chooses per transaction whether to use them.

\medskip

Compression reduces the information disclosed to a recipient but provides no hiding property. Applications must not treat it as a privacy mechanism. Previous holders retain the information they received. The protocols described above provide unlinkability where required.

\medskip

Traceability requirements can be expressed in the token type. A regulated instrument can define transfer rules enforced by every recipient, such as an attestation from an accredited party at each step. Recipients check compliance before acceptance, without requiring retrospective analysis of a public ledger.

\section{History Compression}

A token accumulates one certified transaction per transfer, so an unmodified token grows with use. Compression replaces a prefix of that history with a succinct proof that the prefix was a valid, certified chain.

\medskip

Compression changes the token's representation without executing a transaction, spending state, or creating state. The token identifier, genesis, and current state remain unchanged. The aggregation layer neither participates in nor observes compression. The compressed and original representations identify the same token in the same state, and spending either consumes the same state identifier. Compression requires no token-type opt-in and cannot duplicate a token.

\medskip

The construction uses the append-only property established in Part~\ref{part:protocol}. One recent certificate anchors inclusion of arbitrarily old entries because tree entries are never removed. This replaces a certificate check per transfer with one certificate and a hash path per transfer. Separate proofs of transfer-chain validity and certification allow accumulated history to be compressed to constant size.

\medskip

Compression also applies to the ancestor graph. A token created by a split depends on its source, which may itself have been created by a split. Provenance compression replaces this graph with a certificate about the source alone, allowing one certificate to serve all outputs of a split.

\medskip

Compression is optional and runs in the background, off the payment preparation critical path. It is lossy and irreversible: a holder that discards the original records cannot reconstruct them. This can prevent use by a future counterparty whose verifier lacks the required verification keys. Compression reduces size only above a break-even point; compressing a short history produces a larger token. Wallets should retain raw history separately and use the compressed form for transport.

\section{The Enshrined EVM Partition}\label{sec:evm}

One network partition runs an account-based execution environment compatible with the Ethereum Virtual Machine. The consensus layer certifies its state through the same interface used by aggregation partitions. It hosts governance, staking contracts, and the native currency.

\medskip

The partition is enshrined because its validator set and voting weights derive from the consensus layer's configuration, and execution requires authenticated input from that layer. It is an internal channel within the same system, with no bridge operator. Aggregation partitions operate independently, so EVM delays do not delay certification of unrelated shards.

\begin{figure}[htbp]
    \centering
    \begin{tikzpicture}[diagram]
        \node[diagram box emph, minimum width=11.6cm, minimum height=0.95cm] (core) at (0,2.4)
            {\textbf{Consensus Layer} \quad {\scriptsize orders and certifies}};

        \node[diagram box, minimum width=3.0cm, minimum height=1.5cm, align=center,
              fill=uniFillB, draw=uniAccent] (seal) at (-3.9,0.05)
            {\textbf{Seal}\\\textbf{transaction}\\[2pt]{\scriptsize privileged,}\\{\scriptsize always first}};
        \node[diagram box, minimum width=3.6cm, minimum height=1.5cm, align=center] (user) at (0.2,0.05)
            {\textbf{User transactions}\\[2pt]{\scriptsize ordinary EVM}\\{\scriptsize execution}};
        \node[diagram box, minimum width=3.0cm, minimum height=1.5cm, align=center] (res) at (4.4,0.05)
            {\textbf{New state}\\\textbf{root}\\[2pt]{\scriptsize and block hash}};

        \draw[diagram arrow] (seal) -- (user);
        \draw[diagram arrow] (user) -- (res);

        \draw[diagram arrow] (-3.9,1.925) -- node[diagram label, left, align=right]
            {certified root origin:\\round, time, tree root} (-3.9,0.8);
        \draw[diagram arrow up] (4.4,0.8) -- node[diagram label, right, align=left]
            {certification request} (4.4,1.925);

        \node[diagram box faint, minimum width=11.6cm, minimum height=0.7cm] (note) at (0,-1.5)
            {\scriptsize Contracts read protocol authority only from state written by the seal transaction. Never from a local clock, file, or network response.};
        \draw[diagram connector, draw=uniLine!45] (seal.south) -- (seal.south |- note.north);
        \draw[diagram connector, draw=uniLine!45] (user.south) -- (user.south |- note.north);
    \end{tikzpicture}
    \caption{One EVM block. A privileged system operation imports the certified consensus state before any user transaction runs, and the resulting state root is certified in turn.}
    \label{fig:evm}
\end{figure}

\subsection{The Seal Feed}

Before any user transaction executes, the environment performs exactly one protocol operation that imports the consensus layer's certified state: the round, the epoch, the reference time, the global tree root, and the committed configuration transitions needed to interpret them. Its origin, destination, data, and position are fixed by the protocol, and no private key confers the right to originate it. If the input is missing, forged, stale, or inconsistent, the block is invalid.

\medskip

A contract obtains protocol authority only from state written by this operation. Execution never reads a node-local trust base, clock, file, or network response. Permissionless certificate verification therefore cannot advance the privileged clock or activate a validator set.

\subsection{Partition Functions}

The partition executes staking contracts whose certified output is adopted by the consensus layer. Validators verify the certified contract result without independently rerunning the election algorithm.

\medskip

Contracts can call three consensus primitives at fixed addresses with deterministic cost: verification of a complete Unicity Certificate against authenticated trust-base state, verification of a canonical consensus vote including voting metadata, and verification of an inclusion proof against a supplied authenticated root. These functions allow contracts to verify aggregation-layer state directly. The slashing contract can accept evidence from anyone, and the bridge vault can verify redemptions natively.

\medskip

A bridge vault deployed in this partition authenticates locks through a certificate and a storage proof, without waiting for another chain's probabilistic finality. The two execution models share one root of trust, reducing the work required for transfers between them.

\medskip

EVM-compatible contracts, wallets, and development tools can be used for application components requiring shared state.

\subsection{Native Currency}

UCT is the native currency used for gas, bonded stake, and rewards. Its supply is allocated at genesis, and protocol execution creates no additional currency. Supply therefore equals the genesis allocation less burns. There is no mining, issuance schedule, or earlier supply to migrate.

\medskip

Staking rewards are transfers from an allocated pool rather than new issuance. Slashing transfers collateral to the treasury and neither creates nor destroys currency. Where UCT is represented in another form, as a wrapped contract token for composability or as an execution-layer token claim backed by currency held in a vault, neither representation increases supply, and the accounting keeps the backing from being counted twice.

\medskip

Custody, withdrawal accounting, allocation, and vault contracts are immutable in the initial profile. Stake custody and its slashing and exit rules are separate from election policy. Election policy changes therefore cannot authorize arbitrary withdrawals or slashing. Parameter governance confers no treasury spending authority.

\section{Bridging}\label{sec:bridging}

Bridges are a common source of asset losses, almost always through compromise or action by an operator, committee, or multisig. Unicity's bridge uses none of these in either direction.

\begin{figure}[htbp]
    \centering
    \begin{tikzpicture}[diagram]
        % external chain column
        \node[diagram box ext, minimum width=4.0cm, minimum height=1.15cm, align=center]
            (vault) at (-4.6,0) {\textbf{Vault}\\{\scriptsize pooled custody, external chain}};
        \node[diagram box ext, minimum width=4.0cm, minimum height=0.8cm, align=center]
            (acc) at (-4.6,-1.85) {\textbf{Nullifier root}\\{\scriptsize one value on chain}};
        \draw[diagram connector, draw=uniWarn!70] (vault) -- (acc);

        % unicity column
        \node[diagram box emph, minimum width=4.2cm, minimum height=1.15cm, align=center]
            (tok) at (4.6,0.75) {\textbf{Unicity token}\\{\scriptsize backing proof in its genesis}};
        \node[diagram box, minimum width=4.2cm, minimum height=1.15cm, align=center]
            (burn) at (4.6,-1.85) {\textbf{Burn}\\{\scriptsize commits recipient and amount}};

        % bridge in
        \draw[diagram arrow] (vault.east) -- node[diagram label, above, align=center]
            {\textbf{in:} lock, then mint against\\a proof every recipient checks} (tok.west);
        % bridge out
        \draw[diagram arrow up] (burn.west) -- node[diagram label, below, align=center]
            {\textbf{out:} one succinct proof\\releases a batch of burns} (acc.east);
        \draw[diagram arrow] (tok.south) -- node[diagram label, right, align=left]
            {redeem} (burn.north);

        \node[diagram caption, text width=14.0cm, align=center] at (0,-3.0)
            {The bridge requires no operator, committee, or challenge window. The vault releases funds only against a valid proof.
             The prover can delay release but cannot alter the recipient or amount.};
    \end{tikzpicture}
    \caption{Bridging in both directions. Entry is authorized by a proof that recipients check; exit is authorized by a proof the vault checks for constant cost.}
    \label{fig:bridge}
\end{figure}

\subsection{Bringing an Asset In}

To bridge an asset into Unicity, a party calls the external-chain vault. The vault places the asset in pooled custody, assigns a lock number, and emits an event committing to the representing Unicity token's identifier, the recipient's spending condition, and the amount. The Unicity token is then minted with a reason referencing that lock.

\medskip

Permissionless minting alone establishes no backing. Every recipient runs the type's reason verifier. It confirms that a final lock with the specified number exists for the configured vault, asset, and chain, and that the committed identifier, recipient condition, and amount exactly match the genesis. These checks enforce the following properties:

\begin{itemize}
    \item A token without a matching final lock fails verification and is rejected.
    \item One lock funds exactly one token. The lock commits to the token identifier, which determines the single state at which its genesis is recorded.
    \item The lock commits to the recipient's spending condition. Anyone may construct the genesis, but only the intended recipient can spend the resulting token.
    \item The declared value must equal the locked amount, preventing inflation.
    \item The chain, vault, and asset must match the pinned configuration, excluding unauthorized contracts.
    \item The lock must be final on the external chain, excluding reversed deposits.
\end{itemize}

Accepting an external deposit requires waiting for finality, since a chain reorganization could undo the lock after the token begins circulating. Each verifier enforces this wait individually; it is not a global condition attached to the token. A party minting to a condition it already controls need not wait because it knows whether its own deposit exists. The wait applies to the first party that did not witness the lock.

\subsection{Taking an Asset Out}

A Unicity burn is final upon certification and can be verified offline against the trust base, without further waiting or a node query. Replaying token history is too expensive for the vault contract. A succinct validity proof replaces this history-dependent check with a single constant-cost verification.

\medskip

Replay is prevented by an accumulator whose root is the only bridge state the external chain stores. The nullifier of a burn derives from its certified state, which the aggregation layer records at most once, so it is globally unique by the same mechanism that prevents double-spending. A release proves the nullifier is absent from the accumulator and inserts it, and the vault advances its stored root. The on-chain state required for replay prevention is therefore one value, regardless of how many redemptions have occurred, and anyone can rebuild the accumulator from the external chain's own transaction history.

\medskip

One proof can cover a batch of burns. It establishes that a single recent certificate anchors every involved transition. Each token must pass full verification and end in a burn committing to the stated recipient and amount. Each amount must trace through split conservation to a genesis backed by a lock recorded by the vault. The proof also establishes that each nullifier was absent and has now been inserted.

The vault verifies the proof, checks its configuration and the allow-listed trust base, and requires the accumulator to match its stored root. It confirms each referenced lock against its records, advances the root, and pays. No external query is required because backing was committed locally when the asset was locked.

\subsection{Bridge Security Properties}

Solvency holds by construction. Every circulating bridged token is backed by an equal lock, split verification guarantees children never exceed their parent, and so the sum of valid burns cannot exceed the sum of locks. Over-release would require burning an unbacked token, which the backing check and the vault's own lock records exclude.

\medskip

Double release is excluded because the nullifier is unique, a release both proves its absence and inserts it, and the vault gates on its stored root. A replayed proof fails because the root has moved, and a freshly built one cannot prove absence of something already inserted.

\medskip

Authenticity is enforced inside the proof, which embeds full token verification against the pinned trust base, so a forged or non-final burn is unrepresentable under the proof system's soundness.

\medskip

The proof producer is not trusted. It cannot alter the recipient or amount fixed in the certified burn, or fabricate a release. It can delay proof production, but any party holding the burned token can produce the proof independently.

\medskip

The verification key is immutable after vault deployment, preventing any party, including the deployer, from substituting a different proving relation. Correcting or extending that relation therefore requires a new vault. Value locked in the old vault must be returned under the old relation.
